
\documentclass[10pt]{article}
\usepackage[utf8]{inputenc}
\usepackage[T1]{fontenc}
\usepackage{amsmath, amssymb, xcolor, geometry, enumitem, multicol, tcolorbox}

\definecolor{mainblue}{RGB}{0, 102, 204}
\geometry{a4paper, margin=1.2cm, top=1cm, bottom=3cm, footskip=1.2cm}

\begin{document}
\enlargethispage{2.5cm}

% Modern Header
\begin{tcolorbox}[colback=mainblue!5, colframe=mainblue, sharp corners, boxrule=0.5pt]
    \small \textbf{Unit:} undefined \hfill \textbf{Name:} \rule{3cm}{0.4pt} \\
    \small \textbf{Topic:} Variables and Expressions / Order of Operations \hfill \textbf{Date:} \rule{2cm}{0.4pt} \\
    \centering \textbf{\large Homework 1}
\end{tcolorbox}

\vspace{0.2cm}

% MCQ Section
\begin{multicols}{2}
\begin{enumerate}[label=\textbf{Q\arabic*.}, leftmargin=*, itemsep=6pt, parsep=0pt]

  \item \small What is the variable in the expression $2x$?
  \begin{enumerate}[label=(\Alph*), leftmargin=0.6cm, nosep, itemsep=2pt]
    \item 2
    \item x
    \item 2x
    \item +
    \item None of the above
  \end{enumerate}

  \item \small Identify the constant term in $3y + 2$.
  \begin{enumerate}[label=(\Alph*), leftmargin=0.6cm, nosep, itemsep=2pt]
    \item 3y
    \item 2
    \item +
    \item 3
    \item $3y + 2$
  \end{enumerate}

  \item \small What operation does the symbol $+$ represent in $4 + 5$?
  \begin{enumerate}[label=(\Alph*), leftmargin=0.6cm, nosep, itemsep=2pt]
    \item Subtraction
    \item Addition
    \item Multiplication
    \item Division
    \item None of the above
  \end{enumerate}

  \item \small What is the coefficient of $x$ in $-2x$?
  \begin{enumerate}[label=(\Alph*), leftmargin=0.6cm, nosep, itemsep=2pt]
    \item $-2$
    \item $x$
    \item $2x$
    \item $-x$
    \item $2$
  \end{enumerate}

  \item \small Which part of the expression $4x - 3$ is the term with the variable?
  \begin{enumerate}[label=(\Alph*), leftmargin=0.6cm, nosep, itemsep=2pt]
    \item $4x$
    \item $-3$
    \item $4$
    \item $x$
    \item $-$
  \end{enumerate}

  \item \small What does the term $3 cdot 2$ represent in terms of operations?
  \begin{enumerate}[label=(\Alph*), leftmargin=0.6cm, nosep, itemsep=2pt]
    \item Addition
    \item Subtraction
    \item Multiplication
    \item Division
    \item None of the above
  \end{enumerate}

  \item \small Identify the expression that shows $3$ added to $x$.
  \begin{enumerate}[label=(\Alph*), leftmargin=0.6cm, nosep, itemsep=2pt]
    \item $x + 3$
    \item $x - 3$
    \item $3x$
    \item $3 - x$
    \item $x cdot 3$
  \end{enumerate}

  \item \small What is the order of operations represented by PEMDAS?
  \begin{enumerate}[label=(\Alph*), leftmargin=0.6cm, nosep, itemsep=2pt]
    \item Parentheses, Exponents, Multiplication and Division, and Addition and Subtraction
    \item Parentheses, Multiplication and Division, Exponents, and Addition and Subtraction
    \item Parentheses, Addition and Subtraction, Exponents, and Multiplication and Division
    \item Parentheses, Exponents, Addition and Subtraction, and Multiplication and Division
    \item None of the above
  \end{enumerate}

\end{enumerate}
\end{multicols}

\vspace{-0.2cm}
\hrule
\vspace{0.2cm}
\textbf{\small Open Ended Questions (FRQ)}
\begin{enumerate}[label=\textbf{Q\arabic*.}, start=9, itemsep=5pt, topsep=0pt]

  \item \small Draw a simple algebraic expression that represents the phrase '5 more than $x$'. Explain what each part of your expression represents. \vspace{1.2000000000000002cm}

  \item \small Explain in your own words the difference between a constant and a variable in an algebraic expression. Provide a simple example of each. \vspace{1.2000000000000002cm}

\end{enumerate}

\vfill

% Chef's Tips Footer
\begin{tcolorbox}[colback=blue!5, colframe=mainblue!50, title=\small \textbf{Chef's Tips}, fonttitle=\bfseries, bottom=1mm]
\begin{itemize}[noitemsep, topsep=2pt, leftmargin=0.5cm]
  \item \scriptsize Ensure students understand that variables are letters that represent unknown or changing values.\item \scriptsize Constant terms are numbers without variables, and coefficients are numbers multiplied by variables.\item \scriptsize Students should be able to identify and explain basic operators such as $+$, $-$, $cdot$, and $div$ in simple expressions.
\end{itemize}
\end{tcolorbox}

\end{document}
  